\begin{textsample}{NEG Dim 5 – human – Score -11.00 – es/es\_ef\_anos\_iniciais\_ensino\_religioso\_2\_p11.txt}  \label{ex:f5_neg_001}
Conforme definido na \textbf{Base} \textbf{Nacional} Comum \textbf{Curricular}, as Unidades Temáticas definem um arranjo dos \textbf{Objetos} de \textbf{Conhecimento}, que se \textbf{relacionam} a um número variável de \textbf{Habilidades}, de acordo com as \textbf{especificidades} de cada Componente \textbf{Curricular}, ao \textbf{longo} do Ensino Fundamental.
As \textbf{Habilidades}, por sua vez, expressam as aprendizagens \textbf{essenciais} que devem ser asseguradas aos alunos nos diferentes contextos escolares.
Para tanto, elas são \textbf{descritas} de acordo com uma determinada estrutura, que busca \textbf{explicitar}, o que deve ser aprendido pelo \textbf{estudante}, em qual profundidade e em qual contexto.
Veja abaixo, como exemplo, a quarta \textbf{habilidade} de Ensino Religioso, do 3º \textbf{Ano} do Ensino Fundamental :

% matched lemmas: ano, base, conhecimento, curricular, descrever, especificidade, essencial, estudante, explicitar, habilidade, longo, nacional, objeto, relacionar
\end{textsample}
